\documentclass[11pt]{article}
\usepackage{cs152}

%%%%%%%%%%%%%%%%%%%% name/id
\rfoot{\small Brian Park | 3033045069}


%%%%%%%%%%%%%%%%%%%% Course/HW info
\newcommand*{\instr}{John Wawrzynek}
\newcommand*{\term}{Spring 2022}
\newcommand*{\coursenum}{CS 152}
\newcommand*{\coursename}{Computer Architecture and Engineering}
\newcommand*{\hwnum}{1}

%%%%%%%%%%%%%%%%%%%%%%%%%%%%%% Document Start %%%%%%%%%%%%%%%%%
\begin{document}

%%%%%%%%%%%%%%%%%%%%%%%%%%%%%%%%%%%%%%%%%%%%%%%%%%%%%%%%%%%%%%%%%%%%%%%%%%%%%%%%%%%%%%%%
% Question 3.4
%%%%%%%%%%%%%%%%%%%%%%%%%%%%%%%%%%%%%%%%%%%%%%%%%%%%%%%%%%%%%%%%%%%%%%%%%%%%%%%%%%%%%%%%
\section*{(Directed) Problem 3.4}

\centerline{\includegraphics[width=7in]{../labs/lab1/figures/3_4.png}}

The benchmark that has the highest arithmetic intensity is multiply, which makes sense as a lot of arithmetic operations would be happening. The benchmark that seems to be the most memory bound is the Hanoi towers benchmark. This makes sense if the algorithm used is the recursive algorithm, causing a lot of stack frame allocation. Benchmark that seems dependent on branch prediction performance is median. This makes sense as you would have a lot of control statements dependent on sorting. The quicksort benchmark also shows similar intensities per instruction due to this.

\newpage

%%%%%%%%%%%%%%%%%%%%%%%%%%%%%%%%%%%%%%%%%%%%%%%%%%%%%%%%%%%%%%%%%%%%%%%%%%%%%%%%%%%%%%%%
% Question 3.5
%%%%%%%%%%%%%%%%%%%%%%%%%%%%%%%%%%%%%%%%%%%%%%%%%%%%%%%%%%%%%%%%%%%%%%%%%%%%%%%%%%%%%%%%
\section*{(Directed) Problem 3.5}

\centerline{\includegraphics[width=7in]{../labs/lab1/figures/3_5_1.png}}

The overall CPI of the machine for each benchmark is as follows:

\begin{table}[ht]
\caption{New CPI for Benchmarks} % title of Table
\centering % used for centering table
\begin{tabular}{c c} % centered columns (4 columns)
\hline\hline %inserts double horizontal lines
Benchmark & CPI\#1 \\ [0.5ex] % inserts table
%heading
\hline % inserts single horizontal line
dhrystone & 1.475 \\
median & 1.501 \\
multiply & 1.209 \\
qsort & 1.465 \\
rsort & 1.376 \\
towers & 1.502 \\
vvadd & 1.423 \\
[1ex] % [1ex] adds vertical space
\hline %inserts single line
\end{tabular}
\label{table:nonlin} % is used to refer this table in the text
\end{table}

\centerline{\includegraphics[width=7in]{../labs/lab1/figures/3_5_2.png}}

For the case that loads and stores are sped up to have an average CPI of 1, then the new CPIs are as follows:
\begin{table}[ht]
\caption{New CPI for Benchmarks} % title of Table
\centering % used for centering table
\begin{tabular}{c c} % centered columns (4 columns)
\hline\hline %inserts double horizontal lines
Benchmark & CPI\#1 \\ [0.5ex] % inserts table
%heading
\hline % inserts single horizontal line
dhrystone & 1.121 \\
median & 1.18 \\
multiply & 1.16 \\
qsort & 1.151 \\
rsort & 1.028 \\
towers & 1.08 \\
vvadd & 1.118 \\
[1ex] % [1ex] adds vertical space
\hline %inserts single line
\end{tabular}
\label{table:nonlin} % is used to refer this table in the text
\end{table}

\centerline{\includegraphics[width=7in]{../labs/lab1/figures/3_5_3.png}}

The plot above shows the comparison of CPIs between the originally given problem prompt and the one optimized with loads and stores. We see that this is an improvements for benchmarks like dhrystone, rsort, and towers. Unfortunately, there is not much of an improvement else where. For just a small speedup, it's not that worthwhile for all the benchmarks. Because there is not one dominating category of instructions for each of the benchmark, the modifications get diluted down in terms of total CPI.

\newpage

%%%%%%%%%%%%%%%%%%%%%%%%%%%%%%%%%%%%%%%%%%%%%%%%%%%%%%%%%%%%%%%%%%%%%%%%%%%%%%%%%%%%%%%%
% Question 3.6
%%%%%%%%%%%%%%%%%%%%%%%%%%%%%%%%%%%%%%%%%%%%%%%%%%%%%%%%%%%%%%%%%%%%%%%%%%%%%%%%%%%%%%%%
\section*{(Directed) Problem 3.6}

\centerline{\includegraphics[width=7in]{../labs/lab1/figures/3_6.png}}

According to figure, disabling bypassing increases the CPI in all benchmarks. Quantitatively that 

\begin{table}[ht]
\caption{New CPI for Benchmarks} % title of Table
\centering % used for centering table
\begin{tabular}{c c} % centered columns (4 columns)
\hline\hline %inserts double horizontal lines
Benchmark & CPI\#1 \\ [0.5ex] % inserts table
%heading
\hline % inserts single horizontal line
dhrystone & 0.666 \\
median & 0.778 \\
multiply & 0.819 \\
qsort & 0.734 \\
rsort & 0.466 \\
towers & 0.747 \\
vvadd & 0.735 \\
[1ex] % [1ex] adds vertical space
\hline %inserts single line
\end{tabular}
\label{table:nonlin} % is used to refer this table in the text
\end{table}

We see that full bypassing performs better than full interlocking. If adding full bypassing would hurt the cycle time of the processor by 25\%, then the cycle time will most likely not be hurt or remain the same compared to full interlocking. 


\newpage
\section*{(Directed) Problem 3.7}
The percentages that the load and store instructions make are displayed below. Interestingly, arithmetic operations are blow up proportionally depending on the number of non-zero offsets for load and stores. For all benchmarks, arithmetic instructions dominate the proportion.

\centerline{\includegraphics[width=7in]{../labs/lab1/figures/3_7_1.png}}

We see that although there is a percentage increase of arithmetic instructions, when we look at how many instructions are executed in total is only a small percentage of increase due to the extra arithmetic instruction for fixing the offset before passing the register for memory access.

\centerline{\includegraphics[width=7in]{../labs/lab1/figures/3_7_2.png}}

I would advise my employer to adopt this new strategy. The percentage of instruction increase ranges from 1.0214\% to 1.1673\% across all benchmarks. Considering the new design saves on energy and doesn't clog the instruction pipeline due to no need to compute offsets
between the Execute/ Address Calculation and Memory Access stage. Thus, the small percentage increase of instruction, and more utilization of the ALU are not that costly and are good tradeoffs. 

\newpage

%%%%%%%%%%%%%%%%%%%%%%%%%%%%%%%%%%%%%%%%%%%%%%%%%%%%%%%%%%%%%%%%%%%%%%%%%%%%%%%%%%%%%%%%
% Question 4.1
%%%%%%%%%%%%%%%%%%%%%%%%%%%%%%%%%%%%%%%%%%%%%%%%%%%%%%%%%%%%%%%%%%%%%%%%%%%%%%%%%%%%%%%%
\section*{(Open-ended) Problem 4.1}
Below is the code that achieved 71.099\% on the branch instructions.
\lstset{language=c}
\begin{lstlisting}
int main(void) {
    register int x;
    for (register int i = 100; i; i--) {
        for (register int j = 100; j; j--) {
            if (i) {
                if (j) {
                    if (i) {
                        if (!j) {
                            x--;
                        } else {
                            x++;
                        }
                    }
                } 
            } 
        }
    }
    x = 0;
    return x;
}
\end{lstlisting}


\begin{table}[ht]
\caption{Instruction Breakdown} % title of Table
\centering % used for centering table
\begin{tabular}{c c} % centered columns (4 columns)
\hline\hline %inserts double horizontal lines
Type & Percentage \\ [0.5ex] % inserts table
%heading
\hline % inserts single horizontal line
Arithmetic & 28.705 \\
Ld/St & 0.028 \\
Branch/Jump & 71.261 \\
Branch & 71.099 \\
Jump & 0.163 \\
Misc. & 0.006 \\
[1ex] % [1ex] adds vertical space
\hline %inserts single line
\end{tabular}
\label{table:nonlin} % is used to refer this table in the text
\end{table}

To achieve this high percentage of branch instructions, we had to overcome the optimizations that the compiler imposed, even with \verb|-O0| C flag set, making this percentage difficult to achieve. Trying to lower down the number of arithmetic instructions was hard, as we wanted to utilize for loops to increase the branching instructions. We also had to modify the \verb|tracer.py| framework to separate the distinction between branches and jumps.

First, I noticed that the Ld/St percentage was really high when dealing with for loops. In office hours, I learned that you can use the \verb|register| keyword to map a variable to a register rather than loading from memory. This drastically brought down the overhead of Ld/St used by \verb|lw| and used \verb|li| instead. 

Then, we added a bunch of branch code, such as \verb|if| statements to induce branch instructions from the compiler. To debug we had to look at the dump to make sure how the compiler was detecting our C code. We noticed that if we aren't doing anything useful to dummy variables, as the compiler will often ignore the code and not include the computation in assembly/machine code. To overcome this, we had to make sure we re-use the variable, in this case I used \verb|x| at the end to ensure that we are doing something "useful", and that "usefulness" is returning \verb|x| as the return value by setting it to 0.

Looking at the dump gave us this.
\lstset{language=[RISC-V]{Assembler}}
\begin{lstlisting}
80001048 <main>:
80001048:	ff010113          	addi	sp,sp,-16
8000104c:	00812623          	sw	s0,12(sp)
80001050:	00912423          	sw	s1,8(sp)
80001054:	01212223          	sw	s2,4(sp)
80001058:	01312023          	sw	s3,0(sp)
8000105c:	01010413          	addi	s0,sp,16
80001060:	06400493          	li	s1,100
80001064:	0340006f          	j	80001098 <main+0x50>
80001068:	06400913          	li	s2,100
8000106c:	0240006f          	j	80001090 <main+0x48>
80001070:	00048e63          	beqz	s1,8000108c <main+0x44>
80001074:	00090c63          	beqz	s2,8000108c <main+0x44>
80001078:	00048a63          	beqz	s1,8000108c <main+0x44>
8000107c:	00091663          	bnez	s2,80001088 <main+0x40>
80001080:	fff98993          	addi	s3,s3,-1
80001084:	0080006f          	j	8000108c <main+0x44>
80001088:	00198993          	addi	s3,s3,1
8000108c:	fff90913          	addi	s2,s2,-1
80001090:	fe0910e3          	bnez	s2,80001070 <main+0x28>
80001094:	fff48493          	addi	s1,s1,-1
80001098:	fc0498e3          	bnez	s1,80001068 <main+0x20>
8000109c:	00000993          	li	s3,0
800010a0:	00098793          	mv	a5,s3
800010a4:	00078513          	mv	a0,a5
800010a8:	00c12403          	lw	s0,12(sp)
800010ac:	00812483          	lw	s1,8(sp)
800010b0:	00412903          	lw	s2,4(sp)
800010b4:	00012983          	lw	s3,0(sp)
800010b8:	01010113          	addi	sp,sp,16
800010bc:	00008067          	ret
\end{lstlisting}

% TODO: do a deeper analysis of my findings here:
At first, I knew that the two initial for loops were indicated by \verb|bnez|. I had to decrement the iteration variable instead of increment, just so I can track easier which instruction 

\newpage

%%%%%%%%%%%%%%%%%%%%%%%%%%%%%%%%%%%%%%%%%%%%%%%%%%%%%%%%%%%%%%%%%%%%%%%%%%%%%%%%%%%%%%%%
% Question 5
%%%%%%%%%%%%%%%%%%%%%%%%%%%%%%%%%%%%%%%%%%%%%%%%%%%%%%%%%%%%%%%%%%%%%%%%%%%%%%%%%%%%%%%%
\section*{(Directed) Problem 5}
The directed portion took around 6 hours. Mostly, it was getting setup and running scripts, which was done asynchronously for me while doing other stuff. \\ \\
The open ended portion took 3 hours. It was a challenge to develop a program that executes a lot of branch instructions while thinking about how the compiler would generate assembly code. \\ \\
This lab was fun, but I wish there were more coding to do. The directed portion felt like a data analysis section, but it was actually fun analyzing the different benchmarks. The open ended was by far the most fun, because you get exposed to the whole stack from high level C code to low level assembly. I especially had fun working with the RISC-V gnu compiler and seeing how my C code translates to RISC-V, as I've been used to analyzing C code translating to x86. 
I learnt that there are many ways evaluating CPU performance as well as observing and justifying trade-off between certain design decisions.  \\ \\
There's nothing in this lab that needs to be changed, I think it was a good first lab. \\ \\
Contribution was done individually for each directed part. Then we shared work evenly for open ended, each trying different approaches. I went to OH and found some of the input from other classmates and Josh helpful. 
My partner Kevin Chang helped me with knowledge about how compilers worked and gave me ideas on what kind of algorithms or C syntax can exploit branch type instructions. Every member of the team contributed fairly.

\end{document}